\subsection{The exact triple question and the checked dependencies}

\begin{unresolved}[Part B.1: not settled in this record]\label{B:triple}
We have not proved or refuted that every total category supplies a simultaneous affine equalizer for every three ordinary parallel arrows. Precisely, for $a_0,a_1,a_2:p\to q$, the unproved vanishing is
\[
 T_1\!\left(\coker\bigl(P_q\longrightarrow P_p^2,
 t\longmapsto(t(a_1-a_0),t(a_2-a_0))\bigr)\right)=0.
\]
The propositions below identify this group and check all five proposed algebraic leads. They do not assert its vanishing.
\end{unresolved}

We first record the dependencies actually checked. We read TOTEX A.1--A.5, B.1--B.7 and D.13, the corresponding proofs in TOTASPH C.3--C.6 and C.11--C.13, and the affine-criterion proof in TOTASPH D. We also opened the supplied \citet{IsbellMitchell1976}, in particular Lemma~2.2 and Proposition~2.3 on p.~293, Theorem~3.1 on pp.~294--295, Proposition~3.3 on p.~295 and the coefficient remark on p.~296.

For R1 the step used here is literal: the degree-$n$ generators of the bar complex of $\Z[H]$ consist of a composable $n$-chain in $A$ and an element of $H$ at its source. These are precisely the nerve simplices of $\El_A(H)$, with the same faces. Thus totality makes $T_i(\Z[H])$ vanish for $i>0$, without assuming that $H$ is connected, nonempty, or has injective transition maps. The source's dependent-sum construction correctly transfers actions on an element component to actions on $A$; its cyclic-cover argument applies to unrestricted actions. We do not replace $\Z[H]$ by an arbitrary module.

For R2, Isbell--Mitchell's Theorem~3.1 equates exactness with filtered components of $\Aff_\Z(A)$, and Proposition~3.3 equates this with I and III. The finite-support step in the necessity proof really uses common successors to put a finite sum at one target. The sufficient direction follows directly because affine extension preserves the colimit: for $\sum n_i=1$,
\[
 x-\Bigl(\sum_i n_if_i\Bigr)x=\sum_i n_i(x-f_ix).
\]
For R4, the common-successor argument applies also to the total element category $p/A$, and therefore yields I. The pair presentation and $H_1$ vanishing yield an affine equalizer after finitely many augmentation witnesses have been moved to one target. We give below a direct proof of the instance of Proposition~2.3 that we use. For R7, we have checked the two short exact sequences involving the image of the relation map and the fact that the induced map on representable colimits is zero. In particular its row kernel in $P_q^2$ is kept distinct from the column kernel in $P_q$.

All modules in this section are covariant functors $A\to\Ab$ and every displayed module map is a natural $\Z$-linear transformation. We use the canonical identification $\colim_A P_x=\Z$ by augmentation. For a submodule $E\subset P_x$, $\eps(E)$ means the image of $\colim E\to\Z$, not a pointwise union. When I holds, every finite collection of objects in one component has a common successor: completion of spans shortens a finite connecting zigzag, and induction handles finitely many objects.

\subsection{The intersection and one-relation tests}

\begin{theorem}[Proved: Lead 1]\label{B:lead1}
Let $A$ be total, let $a_0,a_1,a_2:p\to q$, and put
\[
 z_i=a_i-a_0,\quad K_i=\ker(P_q\xrightarrow{\,z_i\,}P_p),\quad
 L=K_1\cap K_2,\quad S=K_1+K_2,\quad C=P_q/S.
\]
Then the surjection of R7 is canonically an isomorphism
\[
 T_1(C)\cong\Z/\eps(L).
\]
Moreover, with $J_i=P_qz_i\subset P_p$ and
$D=\im(P_q\to J_1\oplus J_2,\ t\mapsto(tz_1,tz_2))$,
\[
 D\cong P_q/L,\qquad C\cong(J_1\oplus J_2)/D,
 \qquad \colim J_i=0,\quad T_j(J_i)=0\quad(j>0).
\]
\end{theorem}
\begin{proof}
Define $H_i=A(p,-)/(ca_0\sim ca_i)$. Free abelianization preserves this set coequalizer, giving
\[
 0\longrightarrow J_i\longrightarrow P_p\longrightarrow\Z[H_i]
 \longrightarrow0.
\]
The last module is acyclic by totality. Its colimit is $\Z$: every element of $H_i$ is reached from the distinguished element at $p$, so its category of elements is connected. The induced map $\colim P_p\to\colim\Z[H_i]$ is the identity of $\Z$. The long exact sequence therefore gives $\colim J_i=0$ and $T_j(J_i)=0$ for $j>0$. Applying it to
$0\to K_i\to P_q\to J_i\to0$ proves that the augmentation
$\colim K_i\to\Z$ is an isomorphism and that each $K_i$ is acyclic.

The pointwise exact sequence
\[
 0\longrightarrow L\xrightarrow{\,l\mapsto(l,-l)\,}
 K_1\oplus K_2\xrightarrow{\,(u,v)\mapsto u+v\,}S
 \longrightarrow0
\]
gives
\[
 \colim S\cong\Z^2/\{(n,-n):n\in\eps(L)\}.
\]
The map to $\colim P_q=\Z$ is induced by $(a,b)\mapsto a+b$. Its kernel is canonically $\Z/\eps(L)$ through $(n,-n)\mapsto n$. Since $P_q$ is projective, $0\to S\to P_q\to C\to0$ identifies this kernel with $T_1(C)$. This is precisely the map constructed in R7.

The map defining $D$ has kernel $L$, so $D\cong P_q/L$. The other isomorphism is explicit:
\[
 P_q/(K_1+K_2)\longrightarrow(J_1\oplus J_2)/D,
 \qquad [t]\longmapsto[(tz_1,0)].
\]
It is surjective because subtracting $(vz_1,vz_2)$ from $(uz_1,vz_2)$ leaves $((u-v)z_1,0)$. Its kernel is $K_1+K_2$: the equality $(tz_1,0)=(wz_1,wz_2)$ says $w\in K_2$ and $t-w\in K_1$, and conversely these conditions give that equality. All maps are natural in the target object. Finally $0\to D\to J_1\oplus J_2\to C\to0$ gives the alternative identification $T_1(C)\cong\colim D=\Z/\eps(L)$.
\end{proof}

\begin{theorem}[Proved: Lead 2, including all component and empty cases]\label{B:lead2}
Let $A$ be a small category satisfying I. Its abelian colimit functor is exact if and only if, for all $p,r\in A$ and every $z\in\Z[A(p,r)]$ with $\aug z=0$, there are $s\in A$ and $t\in\Z[A(r,s)]$ such that $\aug t=1$ and $tz=0$. Without any hypothesis on $A$, putting $E=\operatorname{ann}(z)$ gives the canonical formula
\[
 T_1(P_p/P_rz)\cong\Z/\eps(E).
\]
For a total category this annihilator property holds in particular for every binomial $z=c-c'$.
\end{theorem}
\begin{proof}
Let $J=P_rz\subset P_p$. From
$0\to E\to P_r\to J\to0$ we obtain
$\colim J=\Z/\eps(E)$. The map $\colim J\to\colim P_p=\Z$ is zero: its composite with the epimorphism $\Z=\colim P_r\to\colim J$ is multiplication by $\aug z=0$. The sequence $0\to J\to P_p\to P_p/J\to0$, with projective middle term, proves the displayed formula.

If colimits are exact, this Tor group vanishes, so $\eps(E)=\Z$. Represent $1$ by a finite sum of elements of $E$ at possibly different objects. All those objects lie in the component of $r$, and I gives a common successor. Transporting and adding the representatives gives one $t$ with augmentation $1$ and $tz=0$.

Conversely, given parallel affine arrows $\alpha,\beta:p\to r$, apply the assumed property to $z=\alpha-\beta$. It supplies an affine equalizer for them. Condition I gives common successors in each component of $A$, and $\Aff_\Z(A)$ has the same components. Its components are therefore filtered. Theorem~3.1 of \citet{IsbellMitchell1976} gives exactness, or one can extend every module affinely and compute its unchanged colimit over those filtered components. The binomial assertion is R4's pair criterion.

If $z=0$, then $E=P_r$ and $\id_r$ is an annihilator, agreeing with the formula. When $A(p,r)=\varnothing$ this is the only $z$. For the empty category both quantified conditions are vacuous and the zero colimit functor is exact. No global connectedness assumption has been used.
\end{proof}

A theorem for three ordinary arrows would not by itself prove Theorem~\ref{B:lead2}'s annihilator condition: after an affine equalizer is applied to three arrows, later comparisons are generally between affine arrows. We make no induction from ordinary triples to arbitrary finite families.

\begin{theorem}[Proved: Lead 3 without hypotheses on the category]\label{B:lead3}
Let $A$ be any small category, $k\geq0$ finite, and $a_0,\ldots,a_k:p\to q$. For $z_i=a_i-a_0$, define
\[
 d^T:P_q\longrightarrow P_p^k,\quad t\longmapsto(tz_1,\ldots,tz_k),
 \qquad L=\ker d^T,\quad Q_a=\coker d^T.
\]
There are canonical isomorphisms
\begin{align*}
 T_1(Q_a)&\cong\Z/\eps(L),\\
 T_2(Q_a)&\cong\ker(\colim L\to\Z),\\
 T_i(Q_a)&\cong T_{i-2}(L)\qquad(i\geq3).
\end{align*}
If finite sets of objects in the component of $q$ have common successors, the family has a simultaneous affine equalizer if and only if $T_1(Q_a)=0$. The row map
\[
 d:P_q^k\longrightarrow P_p,\qquad(t_i)\longmapsto\sum_i t_i z_i
\]
has cokernel $\Z[A(p,-)/(ca_0\sim\cdots\sim ca_k)]$. The module $Q_a$ has generators $x_1,\ldots,x_k$ at $p$ and the single module relation $\sum_i z_ix_i=0$ at $q$.
\end{theorem}
\begin{proof}
Put $D=\im d^T$. The first short exact sequence
$0\to L\to P_q\to D\to0$ gives
\[
 \colim D=\Z/\eps(L),\quad
 T_1(D)=\ker(\colim L\to\Z),\quad
 T_j(D)=T_{j-1}(L)\quad(j\geq2).
\]
The map $\colim D\to\colim P_p^k=\Z^k$ is zero because its composite with the epimorphism $\colim P_q\to\colim D$ is the vector of augmentations of the $z_i$, all zero. The long exact sequence of $0\to D\to P_p^k\to Q_a\to0$ gives the three asserted formulas. Canonicity is inherited from the connecting morphisms and the representable augmentations.

An affine equalizer is exactly an element of $L(r)$ with augmentation $1$. Such an element gives $\eps(L)=\Z$. Conversely, common successors allow a finite family of representatives witnessing $1\in\eps(L)$ to be transported and added at one target. This proves the equivalence. For $k=0$, $P_p^0=0$, $L=P_q$ and $Q_a=0$, so all formulas remain valid.

Objectwise, the free abelian group on a quotient set is the quotient of the original free group by the subgroup generated by differences of identified elements. Applying this to the specified relations gives the cokernel of $d$. The generators $x_i$ are the identity elements in the $i$th summands of $P_p^k$, and $\sum_i z_ix_i$ is an element of $(P_p^k)(q)$. This types the last presentation. Under totality the row cokernel is acyclic by R1; transposing this matrix does not identify its cokernel with $Q_a$.
\end{proof}

\subsection{Set-colimit equality, abelian-colimit equality, and the mixed cycle}

\begin{theorem}[Proved: the quantified assertions of Lead 4]\label{B:lead4}
Suppose $A$ is quasi-filtered, $M:A\to\Ab$, and $\alpha,\beta\in M(q)$ have the same class in $\colim_A^{\Set}UM$. Then an affine $t:q\to s$ satisfies $M(t)\alpha=M(t)\beta$.

More generally, assuming only I on $A$, exactness of $\colim_A:\Ab^A\to\Ab$ is equivalent to the same conclusion for every $M,q,\alpha,\beta$ whose classes agree in $\colim_A^{\Ab}M$.
\end{theorem}
\begin{proof}
In the first assertion let $X$ be the connected component subfunctor of $UM$ containing $\alpha$ and $\beta$. Condition I on $A$ reduces the finite zigzag in $X$ to arrows $f,g:q\to r$ with $M(f)\alpha=M(g)\beta$, exactly as in \citet[Lemma~2.2]{IsbellMitchell1976}. Condition IV gives an affine $v:r\to s$ with $vf=vg$. Put $t=vf=vg$. Then
\[
 M(t)\alpha=M(v)M(f)\alpha=M(v)M(g)\beta=M(t)\beta.
\]
This directly checks the instance of their Proposition~2.3 that is required.

For necessity of the second assertion, extend $M$ affinely. By exactness and Theorem~3.1, the components of $\Aff_\Z(A)$ are filtered. The affine-extension comparison identifies the abelian colimits. A filtered colimit of abelian groups has the same underlying set as the filtered colimit of its underlying sets: a finite collection of representatives can be moved to one object, and equality is then witnessed after one transition. Thus the equal images of $\alpha$ and $\beta$ are witnessed by a single affine arrow $t$ out of $q$. Components outside that of $q$ contribute separate direct summands and do not affect this equality.

For sufficiency apply the asserted property to $M=P_p$ and parallel affine $\alpha,\beta:p\to q$. Their images in $\colim P_p=\Z$ are both $1$. We obtain an affine equalizer of every affine pair; I then gives filtered components of $\Aff_\Z(A)$ and hence exactness.
\end{proof}

\begin{proposition}[Proved qualification of Lead 4; stronger readings refuted]\label{B:lead4-qualification}
There is a canonical map
\[
 \theta_M:\colim_A^{\Set}UM\longrightarrow U(\colim_A^{\Ab}M).
\]
For a nonempty connected $A$ with common successors it is surjective. Under I and IV, the main affine-equalization question asks whether all pairs in its fibres admit affine equalizers, since pairs equal already in the source do. This is not an equivalence between exactness and injectivity of $\theta_M$, nor an equivalence between exactness and preservation of finite products by set colimits: both stronger assertions are false.
\end{proposition}
\begin{proof}
Every abelian-colimit element is a finite sum of representatives. Common successors put them at one object, where they can be added, proving surjectivity. The assertion about fibres is exactly Theorem~\ref{B:lead4}. With multiple components surjectivity must instead be asserted componentwise, since a sum may involve several components.

For the first stronger assertion take $A=\Delta_{\mathrm{inj}}$, the countable category of nonempty finite ordinals and increasing injections, and $M=P_{[0]}$. Its exactness is the checked R3. An element of $M([n])$ is an integer vector; an increasing injection inserts zero coordinates, hence preserves the ordered list of its nonzero coefficients. At $[2]$ the vectors
\[
 \alpha=(1,1,-1),\qquad\beta=(1,0,0)
\]
have different such lists and therefore belong to distinct components of $\colim^{\Set}UM$. They both have augmentation $1$, so their abelian-colimit images agree. Indeed the explicit affine arrow
\[
 t=(0,1,2)-(0,1,3)+(0,2,3):[2]\longrightarrow[3]
\]
has augmentation $1$ and sends both vectors to $(1,0,0,0)$. The program \path{code/check_delta_patterns.py} checks this identity and the pattern invariant for all increasing injections into target sizes $3$ through $9$; the zero-insertion proof applies to every target size.

For the second stronger assertion, the cospan category on $[0],[0]$ in $\Delta_{\mathrm{inj}}$ has three components, according as the distinguished points satisfy $<$, $=$ or $>$. Increasing injections preserve these relations; each class is connected by its smallest representative. Hence the diagonal is not final, so set colimits do not preserve finite products, although abelian colimits are exact. The informal final sentence of Lead~4 is valid only in its fibrewise affine-equalization formulation, not as either stronger equivalence.
\end{proof}

\begin{theorem}[Proved: Lead 5, with its coefficient-cycle meaning specified]\label{B:lead5}
Fix $a_0,a_1,a_2:p\to q$, write $z_i=a_i-a_0$, and put
\[
 K=\ker(P_q^2\to P_p,\ (t_1,t_2)\mapsto t_1z_1+t_2z_2),
 \quad K_i=\operatorname{ann}(z_i),\quad
 K_3=\operatorname{ann}(z_2-z_1).
\]
Let $\rho_1,\rho_2,\rho_3\in P_q(s)$ have augmentation $1$ and satisfy
$\rho_1z_1=0$, $\rho_2z_2=0$, $\rho_3(z_2-z_1)=0$. Set
$\kappa=(\rho_1-\rho_3,\rho_3-\rho_2)$.
\begin{enumerate}[label=(\roman*)]
\item We have $\kappa\in K(s)$ and $\eps^2(\kappa)=0$. If $A$ is total, its class in $\colim K$ is zero.
\item If $A$ is connected and satisfies I, $F:A\to\Ab$, and $\xi\in F(s)$, then $[\xi]=0$ in $\colim F$ if and only if there are an arrow $u:s\to y$, finitely many parallel pairs $e_j,e'_j:x_j\rightrightarrows y$, and $k_j\in F(x_j)$ such that
\[
 F(u)\xi=\sum_j\bigl(F(e_j)-F(e'_j)\bigr)k_j.
\]
\item If an affine $t:s\to r$ annihilates $\kappa$, then the common value $t\rho_1=t\rho_2=t\rho_3$ is an affine equalizer of the triple. Such a $t$ exists if the parallel pairs in (ii), applied to $F=K$ and $\xi=\kappa$, have a common affine equalizer. For quasi-filtered $A$, one term in (ii) suffices.
\item The natural map
\[
 \psi:K_1\oplus K_2\oplus K_3\longrightarrow K,
 \qquad(k_1,k_2,k_3)\longmapsto(k_1-k_3,k_2+k_3)
\]
has kernel canonically isomorphic to $L=K_1\cap K_2$. For total $A$, putting $C=K/\im\psi$ gives
\[
 T_1(C)\cong\Z/\eps(L).
\]
At each object $x$, $K(x)$ is the integral cycle group of the bipartite multigraph with vertices $A(q,x)\amalg A(p,x)$ and edges $g\to ga_i$ for $i=0,1,2$; $\im\psi(x)$ is the sum of the cycle groups of the three two-colour subgraphs.
\end{enumerate}
\end{theorem}
\begin{proof}
For (i), direct expansion gives
\[
 (\rho_1-\rho_3)z_1+(\rho_3-\rho_2)z_2
 =\rho_1z_1-\rho_2z_2+\rho_3(z_2-z_1)=0.
\]
Both coordinate augmentations vanish. The checked R7 row calculation, applied to total $A$, makes $\colim K\to\Z^2$ an isomorphism, so $[\kappa]=0$. The two coordinates are restrictions of the representable augmentations; $K$ has not been identified with $L$.

For (ii), in fact common successors for finite sets of objects suffice. The colimit is the quotient of $\bigoplus_x F(x)$ by the subgroup generated by
$i_vF(f)k-i_wk$ for arrows $f:w\to v$. If $[\xi]=0$, choose a finite expression
\[
 i_s\xi=\sum_j\bigl(i_{v_j}F(f_j)k_j-i_{w_j}k_j\bigr).
\]
Absorb integer coefficients and signs into the $k_j$. Choose a common successor $y$ of $s$ and all $v_j,w_j$, and one arrow $h_x:x\to y$ for each distinct object appearing. Apply the homomorphism from this finite direct sum to $F(y)$ whose $x$-component is $F(h_x)$. Then $u=h_s$, $e_j=h_{v_j}f_j$ and $e'_j=h_{w_j}$ give exactly the asserted equation. Conversely every parallel difference has zero colimit class and $[F(u)\xi]=[\xi]$. No transition map of $F$ has been assumed injective.

For (iii), $K(t)\kappa=0$ gives $t\rho_1=t\rho_3=t\rho_2$ in $P_q(r)$. The common value has augmentation $1$ and annihilates both $z_1$ and $z_2$. If $v:y\to r$ is a common affine equalizer of the pairs from (ii), applying $K(v)$ to that equation gives $K(vu)\kappa=0$, and $vu$ is affine. With one term, IV supplies $v$. This last deduction requires the stated pair-equalizer hypothesis; I alone would not supply it.

For (iv), the map lands in $K$ because
$k_1z_1+k_2z_2+k_3(z_2-z_1)=0$. Its kernel consists exactly of triples $(l,-l,l)$ with $l\in L$; membership of $l$ in $K_3$ follows from the two defining annihilations. Thus $l\mapsto(l,-l,l)$ is the specified kernel isomorphism.

Assume totality and put $S=\im\psi$. The pair calculation of Theorem~\ref{B:lead1}, also applied to $(a_1,a_2)$, makes every $K_i$ acyclic with augmentation colimit $\Z$. The row calculation makes $K$ acyclic with colimit $\Z^2$. Right exactness of $0\to L\to\bigoplus_{i=1}^3K_i\to S\to0$ gives
\[
 \colim S\cong\Z^3/\{(n,-n,n):n\in\eps(L)\}.
\]
Its map to $\colim K$ is induced by
$f:\Z^3\to\Z^2$, $f(a,b,c)=(a-c,b+c)$.
This map is surjective and has kernel $\{(n,-n,n):n\in\Z\}$. Since $K$ is acyclic, the long exact sequence of $0\to S\to K\to C\to0$ gives the claimed isomorphism. Under its connecting map, $n\bmod\eps(L)$ goes to the class of $(n,-n,n)$ in $\colim S$.

For the graph assertion, orient every edge from the $A(q,x)$ side to the $A(p,x)$ side. Its degree-one group is $P_q(x)^3$, and its boundary is
\[
 (t_0,t_1,t_2)\longmapsto
 (-t_0-t_1-t_2,\ t_0a_0+t_1a_1+t_2a_2).
\]
The natural isomorphism with its cycle group is
\[
 K(x)\longrightarrow Z_1(G_x;\Z),\qquad
 (t_1,t_2)\longmapsto(-t_1-t_2,t_1,t_2),
\]
with inverse taking the last two coordinates. Restricting to colours $0,1$, to $0,2$, or to $1,2$ gives exactly the three summands in the image of $\psi$. Repeated endpoints retain distinct edges, so the argument includes coincident composites and noninjective transition maps.
\end{proof}

\begin{corollary}[Proved: higher derived groups and the precise mixed-cycle witness]\label{B:mixed-cycle}
Under the totality hypotheses of Theorem~\ref{B:lead5}, its module $C=K/\im\psi$ also satisfies
\[
 T_2(C)\cong\ker(\colim L\to\Z),\qquad
 T_j(C)\cong T_{j-2}(L)\quad(j\geq3).
\]
The class $1\bmod\eps(L)$ in $T_1(C)$ has connecting image $[\kappa]\in\colim(\im\psi)$. If the triple has no affine equalizer, this is a nonzero class. It is represented by a bar $1$-cycle with coefficients in $C$, not by a nonzero ordinary homology class of $BA$.
\end{corollary}
\begin{proof}
Write $S=\im\psi$. Acyclicity of $K_1\oplus K_2\oplus K_3$ gives
$T_1(S)=\ker(\colim L\to\Z^3)=\ker(\colim L\xrightarrow{\eps}\Z)$, because the map is $l\mapsto(\eps l,-\eps l,\eps l)$; it also gives $T_i(S)=T_{i-1}(L)$ for $i\geq2$. Acyclicity of $K$ then gives the asserted formulas. They are the same positive derived-group descriptions as for $Q_a$ in Theorem~\ref{B:lead3}, without asserting $C\cong Q_a$.

The identity $\kappa=\psi(\rho_1,-\rho_2,\rho_3)$ has augmentation vector $(1,-1,1)$, which is the connecting image of $1\bmod\eps(L)$ by the preceding proof. Common successors make absence of an affine equalizer equivalent to $\eps(L)\ne\Z$. The bar complex computes $T_1(C)$, so the nonzero class is represented by a bar $1$-cycle. Lift its finitely many coefficients through the pointwise epimorphism $K\to C$; its bar boundary then lies in $S$ and represents $[\kappa]$. This is the precise coefficient interpretation of the last sentence of Lead~5. An ordinary nonzero class in $H_1(BA;\Z)$ cannot be inferred: totality gives $BA$ weakly contractible.
\end{proof}
